\subsection{Estructuras del scheduler}

Las estructuras en \texttt{riscv/state.h} son un calco directo del TP3:

\begin{codesnippet}
\begin{verbatim}
 sched_miner_state_t sched_miners_state[20];
 uint32_t sched_player_row[2];
 uint32_t sched_miners_left[2];
 uint8_t  sched_last_miner[2];
 uint8_t  sched_this_miner;
\end{verbatim}
\end{codesnippet}

La inicialización (\texttt{sched\_init}) hace exactamente lo mismo que la versión x86: marca los 20 slots como muertos,
reinicia contadores y filas. La única diferencia relevante es que, en este port, no hay una ``GDT de TSS'' que resetear:
los slots son simples índices.

\subsection{\texttt{sched\_nextTask}}

El algoritmo del TP3 (alternar entre jugadores y, si el opuesto no tiene mineros vivos, usar uno del jugador actual; si
ninguno hay, usar idle) se preserva:

\begin{codesnippet}
\begin{verbatim}
 int32_t sched_nextTask() {
     uint32_t actual   = (sched_this_miner < 10) ? PLAYER_A : PLAYER_B;
     uint32_t opuesto  = (actual == PLAYER_A) ? PLAYER_B : PLAYER_A;
     uint16_t ix = 0;

     if (aux_get_next_alive_miner(opuesto, &ix)) {
         sched_this_miner = ix;
         sched_last_miner[opuesto] = ix % 10;
         return ix;
     }
     if (aux_get_next_alive_miner(actual, &ix)) {
         sched_this_miner = ix;
         sched_last_miner[actual] = ix % 10;
         return ix;
     }
     return -1;  // idle
 }
\end{verbatim}
\end{codesnippet}

La diferencia clave con el informe x86 es el \emph{tipo de retorno}: ya no es un selector de GDT, sino el índice 0..19
del minero (o $-1$ para idle). Quien consume ese índice es \texttt{schedule\_next} en \texttt{kernel.cpp}, que convierte
el slot a \texttt{TrapFrame*} y hace \texttt{mmu\_switch\_to\_task(slot)}:

\begin{codesnippet}
\begin{verbatim}
 static TrapFrame* schedule_next() {
     int32_t next = sched_nextTask();
     if (next < 0) return yield_to_idle();
     uint32_t slot = (uint32_t)next;
     init_task_if_needed(slot);
     mmu_switch_to_task(slot);
     g_current_slot = next;
     sched_this_miner = (uint8_t)next;
     return &g_tasks[slot].tf;
 }
\end{verbatim}
\end{codesnippet}

\subsection{Despacho de \texttt{ecall}}

La ``rutina de int 47'' del TP3 se vuelve un \emph{branch} dentro de \texttt{trap\_handle} cuando \texttt{scause}$=$8:

\begin{codesnippet}
\begin{verbatim}
 if (cause_code == 8) {            // ecall desde U-mode
     tf->sepc += 4;                // saltar la propia ecall
     poll_uart_input();
     const uint64_t sys_id = tf->a7;

     if (sys_id == 177788) {       // move(d=a0)
         game_move((e_direction_t)tf->a0);
         return yield_to_idle();   // equiv. del "jmp idle" del TP
     }
     if (sys_id == 310311) {       // take() -> a0
         tf->a0 = (uint64_t)game_take();
         return yield_to_idle();
     }
     if (sys_id == 760279) {       // getId() -> a0
         tf->a0 = (uint64_t)game_get_id();
         return tf;                // sin yield
     }
     if (g_current_slot >= 0) sched_kill();
     return yield_to_idle();
 }
\end{verbatim}
\end{codesnippet}

Detalles:

\begin{itemize}
  \item \texttt{sepc += 4}: a diferencia de \texttt{int N} en x86 (que pushea la dirección \emph{siguiente}), en RISC-V
  \texttt{sepc} apunta a la \texttt{ecall} misma. Hay que avanzarlo manualmente al retornar~\cite{riscv-priv}.

  \item \textbf{Yield a idle para \texttt{move}/\texttt{take}}. El TP3 requería que, tras un \texttt{move} o
  \texttt{take}, la tarea Idle ejecutara el resto del quantum. En x86 eso se implementaba con un \texttt{jmp <idle>}
  explícito. En este port se logra devolviendo \texttt{\&g\_idle\_tf} como próximo \texttt{TrapFrame}: \texttt{trap.S}
  lo cargará y \texttt{sret} dejará a idle corriendo hasta el próximo SSIP.

  \item \textbf{\texttt{getId} sin yield}. También fiel al TP3: retorna al mismo minero con \texttt{tf->a0} seteado.

  \item \textbf{Syscall desconocido}. Se mata a la tarea (\texttt{sched\_kill}), tal como en el TP original.
\end{itemize}

\subsection{Intercambio de tareas por cada tick}

El tick es un trap con \texttt{scause}$=$1 (ver sección~\ref{sec:reloj}). La lógica es:

\begin{codesnippet}
\begin{verbatim}
 if (cause_code == 1) {
     asm volatile("csrc sip, %0" : : "r"((uint64_t)(1ull<<1)));
     if (g_spawn_cooldown_a) g_spawn_cooldown_a--;
     if (g_spawn_cooldown_b) g_spawn_cooldown_b--;
     poll_uart_input();
     update_spinner();
     return schedule_next();
 }
\end{verbatim}
\end{codesnippet}

\texttt{schedule\_next} elige la próxima tarea, \texttt{trap\_handle} retorna el \texttt{TrapFrame*} correspondiente y
\texttt{trap.S} se encarga de restaurarlo y saltar. Todo el ``cambio de tarea'' del TP3 (que en x86 era un \texttt{jmp
<sel>:<off>}) termina siendo una \texttt{sret} con un \texttt{TrapFrame} distinto. No es particularmente elegante, pero
funciona.

\subsection{Atención de excepciones y desalojo}

Cualquier excepción distinta de \texttt{ecall} mata a la tarea actual (si había una corriendo) y cede a idle:

\begin{codesnippet}
\begin{verbatim}
 // Any other exception: kill the current miner (if any) and keep going.
 if (g_current_slot >= 0) {
     print(kHeaderExcPrefix, kHeaderExcX, kHeaderY,
           C_BG_BLACK | C_FG_LIGHT_RED);
     print_hex((uint32_t)cause_code, 2,
               kHeaderExcHexX, kHeaderY,
               C_BG_BLACK | C_FG_LIGHT_RED);
     sched_kill();
 }
 return yield_to_idle();
\end{verbatim}
\end{codesnippet}

Esto cubre page faults (\texttt{scause} 12/13/15), illegal instructions (\texttt{scause}$=$2), misaligned loads/stores,
etc. Por diseño, \texttt{scause} se imprime en hex en el banner superior: es el análogo reducido del ``modo debug'' del
TP3.

\subsection{Modo debug}

El TP3 tenía un modo debug activable con la tecla \texttt{y} que mostraba el estado completo del procesador. En este
port, la información se redujo a un indicador mínimo en la fila 0 del ``video'' (\texttt{TASK EXC XX}) con los dos
dígitos hex del \texttt{cause\_code}. El estado completo de registros no se imprime en pantalla; para debug profundo se
usa QEMU (\texttt{-d int,guest\_errors,cpu\_reset,mmu}). Esta diferencia está documentada en la
sección~\ref{sec:limitaciones}.
